% grf_essay_2026.tex
% Essay for the Gravity Research Foundation 2026 Awards
% "The Price of Time Travel: Three Information-Theoretic Barriers
%  from Minimum Fisher Information"
% Author: Sheng-Kai Huang
% Compile: pdflatex grf_essay_2026.tex && pdflatex grf_essay_2026.tex

\documentclass[12pt]{article}

\usepackage[margin=1in]{geometry}
\usepackage{amsmath}
\usepackage{amssymb}
\usepackage{hyperref}
\usepackage{graphicx}
\usepackage{booktabs}
\usepackage{array}

% Double spacing
\renewcommand{\baselinestretch}{2.0}
\selectfont

% Page numbers
\pagestyle{plain}

% Minimal macros
\newcommand{\rs}{r_{\mathrm{s}}}

\begin{document}

% ============================================================
% TITLE PAGE
% ============================================================
\begin{titlepage}
\renewcommand{\baselinestretch}{1.0}\selectfont
\begin{center}

\vspace*{1cm}

{\LARGE \textbf{The Price of Time Travel:\\[6pt]
Three Information-Theoretic Barriers\\[6pt]
from Minimum Fisher Information}}

\vspace{1.2cm}

{\large Sheng-Kai Huang}

\vspace{0.4cm}

Independent Researcher\\
akai@fawstudio.com\\
Taipei, Taiwan

\vspace{0.8cm}

March 29, 2026

\vspace{1.2cm}

\begin{minipage}{0.85\textwidth}
\textbf{Abstract.}
We show that the gravitational entropy production
$\Sigma = -\ln(-g_{00})$, derived from channel extensivity,
satisfies the vacuum equation $\nabla^2\Sigma = 0$.
The maximum principle for this harmonic field establishes three
barriers to exotic spacetime structures, each with a quantifiable
``price'' to overcome:
(1)~event horizons require stellar-mass energy concentration---a price
Nature already pays;
(2)~warp bubbles require $E \propto v^4 R^2/\delta$ of localized
energy---enormously expensive but finite;
(3)~closed timelike curves require $\Sigma < 0$, meaning regions where
time runs faster than at infinity---the most expensive barrier, requiring
physics beyond the harmonic approximation.
We compute each price explicitly and identify experimental signatures
at the boundary of Barrier~1 testable by the Event Horizon Telescope.
\end{minipage}

\vspace{1.2cm}

Essay written for the Gravity Research Foundation 2026 Awards for Essays on Gravitation.

\end{center}
\renewcommand{\baselinestretch}{2.0}\selectfont
\end{titlepage}

% ============================================================
% PAGE 1-2: THE IDEA
% ============================================================
\setcounter{page}{1}

\section*{I.\ \ The Idea}

How much does a time machine cost?

Not in dollars---in physics. Every science-fiction device that
manipulates spacetime---warp drives, wormholes, time machines---violates
some property of the gravitational field. But violations have prices.
This essay calculates them.




The starting point is an observation Einstein made in 1911, four years
before general relativity \cite{Einstein1911}: the coordinate speed of
light slows down near a massive body, $c_{\mathrm{eff}} = c\sqrt{-g_{00}}$.
Dicke reformulated this as a refractive-medium picture in 1957
\cite{Dicke1957}. We define
\begin{equation}
\Sigma \;=\; -\ln(-g_{00})\,,
\label{eq:sigma}
\end{equation}
and call it the \emph{information cost} of gravity. This is not an ad
hoc definition. In quantum information theory, the gravitational redshift
acts as a pure-loss bosonic channel with transmissivity
$\eta = -g_{00}$ \cite{AhmadiFuentes2014}. If a photon traverses two
gravitational regions in succession, the transmissivities multiply:
$\eta_{\mathrm{tot}} = \eta_1 \cdot \eta_2$. Any extensive measure of
information loss must then satisfy $\Sigma(\eta_1 \eta_2) = \Sigma(\eta_1)
+ \Sigma(\eta_2)$---the Cauchy functional equation, whose unique
continuous solution is $\Sigma = -\ln\eta$ \cite{Huang2026}. Channel
extensivity \emph{forces} equation~(\ref{eq:sigma}).

What field equation does $\Sigma$ satisfy in vacuum? The principle of
minimum information gradient says: minimize the Fisher information---the Cram\'er--Rao lower bound on how precisely the mass distribution can be inferred from the field---
$\mathcal{F}[\Sigma] = \int|\nabla\Sigma|^2\,d^3x$. The
Euler--Lagrange equation is the Laplace equation, $\nabla^2\Sigma = 0$.
The spherically symmetric solution with $\Sigma \to 0$ at infinity and
the correct Newtonian limit is $\Sigma = \rs/r$, giving the exponential
metric $g_{00} = -e^{-\rs/r}$ \cite{Papapetrou1954,Yilmaz1958}.

We emphasize that $\nabla^2\Sigma = 0$ is \emph{not} a consequence of
Einstein's equations---it is an \emph{alternative} vacuum principle.
Einstein's equations in vacuum ($R_{\mu\nu} = 0$) yield the
Schwarzschild metric; Fisher information minimization yields the
exponential metric.  Both agree through second post-Newtonian order
($\beta = \gamma = 1$), diverging only at $O((\rs/r)^3) \sim 10^{-18}$
for solar-system scales.  Which principle Nature selects is an
empirical question---one that the Event Horizon Telescope can answer.

Here is the key insight. The \emph{maximum principle} for harmonic
functions---that $\Sigma$ cannot attain a local extremum in the interior
of the domain---simultaneously forbids three exotic spacetime structures.
Each forbidden structure can be ``purchased'' by supplying enough energy
to make $\nabla^2\Sigma \neq 0$, and the price is calculable.



Gravity, viewed through the $\Sigma$ lens, is spacetime's price list.
The maximum principle is the budget constraint. This essay reads the
receipt.


% ============================================================
% PAGE 2-3: THE MAXIMUM PRINCIPLE
% ============================================================
\section*{II.\ \ The Maximum Principle}

The engine of everything that follows is a classical theorem from the
theory of partial differential equations \cite{Evans2010}. We state it
in the form we need:




\noindent\textbf{Theorem} (Strong Maximum Principle).
\textit{Let\/ $\nabla^2\Sigma = 0$ in a connected domain\/ $D$, with\/
$\Sigma \to 0$ at infinity. Then:}

\begin{itemize}
\item[(i)] $\Sigma > 0$ \textit{everywhere in\/ $D$
(no negative values);}
\item[(ii)] $\sup(\Sigma)$ \textit{is attained only on\/
$\partial D$ (no interior maximum);}
\item[(iii)] $\Sigma$ \textit{is uniquely determined by boundary
conditions.}
\end{itemize}



The proof fits in two lines of any PDE textbook. But its consequences
for gravity are profound. Property~(i) says $g_{00} = -e^{-\Sigma}$
never crosses $-1$: time always runs \emph{slower} than at infinity.
Property~(ii) says $\Sigma$ cannot pile up at a point: no infinite
redshift surface. Property~(iii) says the gravitational field is rigid,
completely fixed by what happens at the boundary.

Gravity, rewritten in the $\Sigma$ language, inherits all the rigidity
of harmonic analysis. The three barriers to exotic spacetimes are three
corollaries of (i)--(iii). And every violation has a price, because the
only way past a harmonic barrier is to break harmonicity---which costs
energy.


% ============================================================
% PAGE 3-4: BARRIER 1 --- EVENT HORIZONS
% ============================================================
\section*{III.\ \ Barrier 1: Event Horizons}

An event horizon requires $g_{00} \to 0$ at some finite radius, which
means $\Sigma \to \infty$. But property~(ii) of the maximum principle
forbids any interior extremum of a harmonic function, let alone a
divergence. The harmonic solution $\Sigma = \rs/r$ gives
$g_{00} = -e^{-\rs/r}$, which is strictly negative for all $r > 0$.
There is no horizon.

To \emph{get} a horizon, one must break harmonicity:
$\nabla^2\Sigma \neq 0$, which requires a source. In the presence of
matter, $\Sigma$ satisfies a Poisson equation
$\nabla^2\Sigma = 8\pi G\rho/c^2$, and concentrating a mass $M$ into a
radius $r \leq \rs = 2GM/c^2$ creates a region where $\Sigma$ is driven
past all bounds by the matter source. The Schwarzschild metric emerges
when $\rho$ is compressed to a point.




\noindent\textbf{The price of Barrier~1:} concentrate mass $M$ into
$r \leq 2GM/c^2$.



Nature already pays this price---but only inside matter.
The vacuum field equation $\nabla^2\Sigma = 0$ applies outside
the star; inside, the Poisson equation $\nabla^2\Sigma = 8\pi G\rho/c^2$
takes over. The horizon forms at the boundary of the matter
distribution during collapse, where $\Sigma$ transitions from the
Poisson regime to the harmonic regime.  The open question is:
what is the vacuum extension beyond the collapsing matter?
Schwarzschild (with horizon) or exponential (without)?
The 4.6\% shadow difference is the empirical arbiter.

Every stellar-mass black hole is a
purchased violation of the maximum principle, financed by
gravitational collapse. Every supermassive black hole at a galactic
center is a receipt for that transaction.

But here is the surprise. The ``cheapest'' solution---the one with
$\nabla^2\Sigma = 0$ everywhere, no source, no horizon---predicts a
\emph{different} photon sphere. In isotropic coordinates, the
exponential metric's effective potential for photons,
$V_{\mathrm{eff}} = e^{-2\rs/\rho}/\rho^2$, has its maximum at
$\rho_{\mathrm{ph}} = \rs$, giving a critical impact parameter
$b_c = e \cdot \rs \approx 2.72\,\rs$, compared with the Schwarzschild
value $b_c = 3\sqrt{3}\,\rs/2 \approx 2.60\,\rs$. The difference is
4.6\%.  (We note that the Kretschmann scalar
$R_{\mu\nu\rho\sigma}R^{\mu\nu\rho\sigma} = \rs^4 e^{-2\rs/r}(3\rs^2
+ 12\rs r + 20 r^2)/(4r^{10})$ vanishes as $r \to 0$; the exponential
metric is singularity-free.)

The Event Horizon Telescope's measurement of the M87* shadow currently
has ${\sim}10\%$ uncertainty \cite{EHT2019}. The next-generation EHT,
with additional baselines and shorter wavelengths, aims for ${\sim}1\%$
precision. At that level, Barrier~1 becomes directly testable: does
Nature pay the \emph{full} price (Schwarzschild, with horizon), or only
the \emph{minimum} price (exponential, without)?


% ============================================================
% PAGE 4-6: BARRIER 2 --- WARP DRIVES
% ============================================================
\section*{IV.\ \ Barrier 2: Warp Bubbles}

A warp bubble is a region of modified spacetime that moves at velocity
$v$ while the interior remains flat. In the $\Sigma$ language, this
requires $\Sigma = 0$ in a flat interior, $\Sigma > 0$ in a thin shell
of thickness $\delta$, and $\Sigma = 0$ in the exterior. That is: a
local maximum of $\Sigma$ confined to a shell of radius $R$.

Property~(ii) of the maximum principle forbids this. If $\Sigma$ is
harmonic, it cannot have a local maximum in the interior.

To purchase a warp bubble, one must again break harmonicity. The
required energy density in the shell sources $\nabla^2\Sigma \neq 0$
and can be estimated by dimensional analysis from the field equation.
The Alcubierre warp metric \cite{Alcubierre1994} in the lapse sector
gives a shell amplitude $\Sigma_0 \sim v^2/c^2$, and the integrated
energy is
\begin{equation}
E_{\mathrm{warp}} \;\sim\;
\frac{c^4}{8G}\,\frac{v^4}{c^4}\,\frac{R^2}{\delta}\,,
\label{eq:warp}
\end{equation}
where $R$ is the bubble radius and $\delta$ the wall thickness.

The $v^4$ scaling is the key. Subluminal warp is exponentially cheaper
than superluminal warp.

\begin{table}[h]
\renewcommand{\baselinestretch}{1.2}\selectfont
\centering
\begin{tabular}{@{}l c l@{}}
\toprule
Speed & Energy & Mass equivalent \\
\midrule
$c$ (light speed)  & $\sim 10^{28}$\,kg  & 3 Jupiter masses \\
$0.01\,c$          & $\sim 10^{20}$\,kg  & large asteroid \\
$100$\,m/s         & $\sim 67$\,kg        & one human \\
$1$\,m/s           & $\sim 60$\,MJ        & 1.5 liters of gasoline \\
\bottomrule
\end{tabular}
\caption{Price list for a warp bubble of radius $R = 10$\,m and wall
thickness $\delta = 1$\,m. The $v^4$ scaling makes subluminal warp
extraordinarily cheap---in principle.}
\label{tab:warp}
\end{table}

The price of a walking-speed warp bubble is 1.5~liters of gasoline.
The catch: in the lapse formulation, it must be \emph{negative}
energy, shaped into a specific $\Sigma$ profile. Nobody knows how
to do this.

A crucial caveat: Fuchs \emph{et al.}\ \cite{Fuchs2024} have shown
that subluminal warp does \emph{not} require negative energy when
formulated via the shift vector rather than the lapse. Our $\Sigma$
framework is purely a lapse-sector analysis, so it overestimates the
difficulty: the true price may be lower than
equation~(\ref{eq:warp}) suggests. Whether the lapse and shift
formulations can be reconciled within the information-theoretic
framework is an open question.

What is \emph{not} an open question is the scaling. Whatever the sign
of the required energy, the $v^4$ dependence means that near-light-speed
travel is vastly more expensive than gentle subluminal nudges. The
universe's price list has a steep luxury tax on speed.


% ============================================================
% PAGE 6-8: BARRIER 3 --- CLOSED TIMELIKE CURVES
% ============================================================
\section*{V.\ \ Barrier 3: Closed Timelike Curves}

Now we come to time travel itself.

A closed timelike curve (CTC) requires, among other conditions, a
region where $g_{00} < -1$, which means $\Sigma < 0$: time runs
\emph{faster} than at infinity. This is a \emph{necessary} condition
for CTCs---the spacetime must also have the right causal topology to
close a timelike path---but the maximum principle forbids even this
first step. Property~(i) states that $\Sigma > 0$ everywhere for a
harmonic field with $\Sigma \to 0$ at infinity. Without $\Sigma < 0$,
no CTC can form. Barrier~3 is qualitatively different from Barriers~1
and~2. Those required large $\Sigma$ (strong gravity); this requires
$\Sigma$ of the \emph{wrong sign}.




\noindent\textbf{The price of Barrier~3:} a ``negative information
source,'' meaning $\nabla^2\Sigma < 0$---a region that \emph{accelerates}
time rather than slowing it.



What does $\Sigma < 0$ mean physically? At $\Sigma > 0$, clocks tick
slower than at infinity (normal gravity; the GPS correction). At
$\Sigma = 0$, flat space, normal time. At $\Sigma < 0$, clocks tick
\emph{faster} than at infinity. Push $\Sigma$ negative enough, past a
threshold, and the light-cone tips over: a CTC forms.

Does $\Sigma < 0$ exist anywhere in known physics?

\emph{The Kerr black hole.}---In the ergosphere of a spinning black
hole, $g_{00} > 0$: the time coordinate becomes spacelike. This is
actually \emph{worse} than $\Sigma < 0$---it is a signature change where
$\Sigma = -\ln(-g_{00})$ is not even defined. The Penrose process
extracts energy from the ergosphere, but does not directly create a CTC.

\emph{The G\"odel universe.}---G\"odel's 1949 rotating cosmology
\cite{Goedel1949} has $g_{00} < -1$ in certain regions, and it
famously contains CTCs. But it requires the \emph{entire universe} to
rotate uniformly. The price is a global constraint: a specific stress-energy tensor ($\rho + p > 0$ with uniform rotation) pervading the entire universe.

\emph{The Alcubierre--CTC connection.}---In special relativity,
faster-than-light travel implies time travel: two superluminal trips
with a Lorentz boost between them close into a CTC. So a superluminal
warp drive (Barrier~2 at $v > c$) automatically triggers Barrier~3.
The two barriers are \emph{independent} for subluminal warp, but
\emph{linked} for superluminal warp. This is why Hawking conjectured
the chronology protection principle \cite{Hawking1992}: Nature may
enforce a cosmic censorship not only of singularities but of time
machines.

In the $\Sigma$ framework, chronology protection is not a conjecture
but a theorem---\emph{within the harmonic approximation}. The maximum
principle \emph{proves} that $\Sigma < 0$ is impossible if
$\nabla^2\Sigma = 0$. The question is whether the universe allows
sources that violate harmonicity in the required direction.




\noindent\textbf{The combined price of a time machine:}
\begin{enumerate}
\item Overcome Barrier~2: supply ${\sim}10^{28}$\,kg of
energy in a thin shell (for light-speed warp).
\item Overcome Barrier~3: create a $\Sigma < 0$ region---currently
unknown physics.
\item These are \emph{independent} barriers. Both must be paid.
\end{enumerate}

A time machine is the most expensive object in the universe. Its price
tag has two parts: an engineering challenge (Barrier~2, enormous but
finite) and a physics challenge (Barrier~3, currently unknown). We have
computed the first; the second remains the province of future theory.


% ============================================================
% PAGE 8-9: THE PRICE LIST
% ============================================================
\section*{VI.\ \ The Price List}

We can now read the universe's receipt. Table~\ref{tab:pricelist}
collects all three barriers, their physical meaning, what breaks them,
and the cost.

\begin{table}[h]
\renewcommand{\baselinestretch}{1.2}\selectfont
\centering
\begin{tabular}{@{}c p{2.2cm} p{2.4cm} p{3.2cm} l@{}}
\toprule
Barrier & Forbidden & What breaks it & Price & Status \\
\midrule
1.\ Horizon &
$\Sigma \to \infty$ at finite $r$ &
Mass concentration &
$M > c^2 r / (2G)$ &
Already paid \\[4pt]
2.\ Warp &
Interior max of $\Sigma$ &
Energy source ($\nabla^2\Sigma \neq 0$) &
$E \sim c^4 v^4 R^2 / (8Gc^4\delta)$ &
Finite but vast \\[4pt]
3.\ CTC &
$\Sigma < 0$ &
Negative info source ($\nabla^2\Sigma < 0$) &
Unknown &
New physics \\
\bottomrule
\end{tabular}
\caption{The price list. All three barriers follow from a single
theorem---the maximum principle for the harmonic information field
$\Sigma$. Gravity is the universe's price list for manipulating time.}
\label{tab:pricelist}
\end{table}

Remarkably, all three prices emerge from a single two-line theorem.
The maximum principle for harmonic functions, applied to the
gravitational information field $\Sigma = -\ln(-g_{00})$, gives the
complete tariff schedule. Property~(i) forbids time machines.
Property~(ii) forbids warp bubbles and horizons. Property~(iii)
guarantees that the cheapest solution---the exponential
metric---is unique.

The exponential metric itself, $g_{00} = -e^{-\rs/r}$, is the
``zero-price'' gravitational field: the solution that satisfies the
vacuum equation everywhere, with minimum Fisher information, and no
exotic structures. It passes all solar-system tests through second
post-Newtonian order ($\beta = \gamma = 1$ in the PPN framework
\cite{Will2014}), with deviations from Schwarzschild appearing only
at $O((\rs/r)^3) \sim 10^{-18}$ for the Sun---far below the Cassini
bound \cite{Bertotti2003}. But in the strong field, near compact
objects, the 4.6\% shadow difference becomes accessible.

The connection to prior work runs deep. Jacobson \cite{Jacobson1995}
showed in 1995 that Einstein's equation is an equation of state;
Dorau and Much \cite{DorauMuch2025} derived it from the quantum
relative entropy in 2025. Our contribution is to show that the
\emph{same} information-theoretic structure that gives the field
equation also gives the \emph{barriers}---the maximum principle is
not an add-on but a built-in feature of the information field.
The barriers are the field equation's way of saying ``no.''


% ============================================================
% PAGE 9-10: OUTLOOK
% ============================================================
\section*{VII.\ \ Outlook}

Three observations close this essay.

\emph{First}, the exponential metric---the harmonic, horizonless,
minimum-Fisher-information solution---is testable now. The ngEHT will
reach the ${\sim}1\%$ shadow precision needed to distinguish
$2.72\,\rs$ from $2.60\,\rs$. If the exponential metric is correct,
there is no event horizon, and the black hole information paradox
takes on a new character---not requiring information to escape from
behind a horizon, because the horizon was never there. The most expensive item on the
price list (Barrier~1) turns out to be the one Nature may not have
purchased after all.

\emph{Second}, the $v^4$ scaling of warp energy makes subluminal warp
tantalisingly cheap in magnitude but frustratingly expensive in kind
(negative energy, in the lapse formulation). The recent positive-energy
warp solutions of Fuchs \emph{et al.}~\cite{Fuchs2024} suggest that
the shift sector may offer a way around the negative-energy
requirement. In the $\Sigma$ framework, this would correspond to a
generalization of the harmonic condition from the lapse to the full
metric---a direction we are actively exploring.

\emph{Third}, the most profound lesson is negative. The maximum
principle does not merely make time travel expensive; it makes it
\emph{categorically} different from the other two barriers. Horizons
and warp bubbles require large $\Sigma$---a quantitative challenge. CTCs
require $\Sigma < 0$---a qualitative one, demanding sources of a
type not known to exist. This is not a no-go theorem---it is a
\emph{price-go} theorem. It says exactly what must be paid, and for
Barrier~3, the currency has not yet been invented.



The question, then, is not whether time travel is possible. The
maximum principle tells us it is forbidden in the harmonic
approximation. The question is: what is the price of going beyond
harmonic? And can the universe afford it?

\bigskip

Gravity is spacetime's price list. The maximum principle is the budget
constraint. And time travel is the item we cannot yet afford.

\bigskip

\noindent\textbf{Acknowledgments.}
The author thanks M.~M.~Wilde for correspondence on the Petz recovery
map, and the gravity and quantum-information communities whose work made
this synthesis possible.


% ============================================================
% REFERENCES
% ============================================================
\newpage
\label{refs_start}
% References start here --- these pages do not count toward the 10-page limit
\begingroup\sloppy
\renewcommand{\baselinestretch}{1.2}\selectfont
\begin{thebibliography}{99}

\bibitem{Einstein1911}
A.~Einstein,
``On the influence of gravitation on the propagation of light,''
Ann.\ Phys.\ (Leipzig) \textbf{35}, 898 (1911).

\bibitem{Dicke1957}
R.~H.~Dicke,
``Gravitation without a principle of equivalence,''
Rev.\ Mod.\ Phys.\ \textbf{29}, 363 (1957).

\bibitem{AhmadiFuentes2014}
M.~Ahmadi, D.~E.~Bruschi, C.~Sab\'{i}n, G.~Adesso, and I.~Fuentes,
``Relativistic quantum metrology: Exploiting relativity to improve quantum
measurement technologies,''
Sci.\ Rep.\ \textbf{4}, 4996 (2014).

\bibitem{Huang2026}
S.-K.~Huang,
``The arrow of time from Petz recovery,''
Zenodo, doi:10.5281/zenodo.14606579 (2026).

\bibitem{Papapetrou1954}
A.~Papapetrou,
``Eine Theorie des Gravitationsfeldes mit einer Feldfunktion,''
Z.\ Phys.\ \textbf{139}, 518 (1954).

\bibitem{Yilmaz1958}
H.~Yilmaz,
``New approach to general relativity,''
Phys.\ Rev.\ \textbf{111}, 1417 (1958).

\bibitem{Evans2010}
L.~C.~Evans,
\emph{Partial Differential Equations}, 2nd ed.
(American Mathematical Society, Providence, 2010).

\bibitem{EHT2019}
Event Horizon Telescope Collaboration,
``First M87 Event Horizon Telescope results. I.\ The shadow of the
supermassive black hole,''
Astrophys.\ J.\ Lett.\ \textbf{875}, L1 (2019).

\bibitem{Alcubierre1994}
M.~Alcubierre,
``The warp drive: hyper-fast travel within general relativity,''
Class.\ Quantum Grav.\ \textbf{11}, L73 (1994).

\bibitem{Fuchs2024}
E.~Fuchs, C.~Cattoen, and M.~Visser,
``Positive-energy subluminal warp drive,''
Class.\ Quantum Grav.\ \textbf{41}, 095001 (2024).

\bibitem{Goedel1949}
K.~G\"odel,
``An example of a new type of cosmological solutions of Einstein's field
equations of gravitation,''
Rev.\ Mod.\ Phys.\ \textbf{21}, 447 (1949).

\bibitem{Hawking1992}
S.~W.~Hawking,
``Chronology protection conjecture,''
Phys.\ Rev.\ D \textbf{46}, 603 (1992).

\bibitem{Kerr1963}
R.~P.~Kerr,
``Gravitational field of a spinning mass as an example of algebraically
special metrics,''
Phys.\ Rev.\ Lett.\ \textbf{11}, 237 (1963).

\bibitem{Penrose1965}
R.~Penrose,
``Gravitational collapse and space-time singularities,''
Phys.\ Rev.\ Lett.\ \textbf{14}, 57 (1965).

\bibitem{DorauMuch2025}
A.~Dorau and A.~Much,
``Quantum relative entropy, Einstein's equations, and spacetime curvature,''
Phys.\ Rev.\ Lett.\ \textbf{134}, 011601 (2025).

\bibitem{Jacobson1995}
T.~Jacobson,
``Thermodynamics of spacetime: The Einstein equation of state,''
Phys.\ Rev.\ Lett.\ \textbf{75}, 1260 (1995).

\bibitem{Will2014}
C.~M.~Will,
``The confrontation between general relativity and experiment,''
Living Rev.\ Relativ.\ \textbf{17}, 4 (2014).

\bibitem{Bertotti2003}
B.~Bertotti, L.~Iess, and P.~Tortora,
``A test of general relativity using radio links with the Cassini spacecraft,''
Nature \textbf{425}, 374 (2003).

\bibitem{Boonserm2018}
P.~Boonserm, T.~Ngampitipan, and M.~Visser,
``Exponential metric represents a traversable wormhole,''
Phys.\ Rev.\ D \textbf{98}, 084048 (2018).

\end{thebibliography}
\endgroup

\end{document}
